\documentclass[10pt,a4paper]{article}

\usepackage[margin=16mm,top=14mm,bottom=15mm]{geometry}
\usepackage[T1]{fontenc}
\usepackage[default]{lato}
\usepackage{xcolor}
\usepackage{microtype}
\usepackage{hyperref}
\usepackage{enumitem}
\usepackage{tabularx}
\usepackage{array}
\usepackage{titlesec}
\usepackage{fancyhdr}
\usepackage{lastpage}
\usepackage{ragged2e}
\usepackage{needspace}
\usepackage{setspace}

\definecolor{Navy}{HTML}{17324D}
\definecolor{Teal}{HTML}{0F7182}
\definecolor{Ink}{HTML}{17202A}
\definecolor{Muted}{HTML}{5D6773}
\definecolor{Rule}{HTML}{D6DEE5}
\definecolor{Soft}{HTML}{EEF4F7}

\hypersetup{
  colorlinks=true,
  urlcolor=Teal,
  linkcolor=Navy,
  pdfauthor={Jia-Hao Ji},
  pdftitle={Jia-Hao Ji - Curriculum Vitae},
  pdfsubject={Biomedical AI, spatial omics, and multi-agent reasoning},
  pdfkeywords={biomedical AI, agentic AI, spatial omics, multi-agent systems, neuroscience}
}
\urlstyle{same}

\pagestyle{fancy}
\fancyhf{}
\renewcommand{\headrulewidth}{0pt}
\renewcommand{\footrulewidth}{0pt}
\fancyfoot[L]{\footnotesize\color{Muted}Jia-Hao Ji \textbar{} Curriculum Vitae \textbar{} September 2026}
\fancyfoot[R]{\footnotesize\color{Muted}Page \thepage{} of \pageref*{LastPage}}
\setlength{\footskip}{9mm}

\setlength{\parindent}{0pt}
\setlength{\parskip}{0pt}
\setstretch{1.03}
\emergencystretch=1em

\titleformat{\section}
  {\needspace{3\baselineskip}\large\bfseries\color{Navy}}
  {}{0pt}{}
  [\vspace{-0.35em}\color{Rule}\titlerule]
\titlespacing*{\section}{0pt}{1.05em}{0.55em}

\newcommand{\contactsep}{\hspace{0.5em}\textcolor{Rule}{\textbullet}\hspace{0.5em}}
\newcommand{\roleline}[3]{%
  \begin{tabularx}{\textwidth}{@{}>{\RaggedRight\arraybackslash}X>{\RaggedLeft\arraybackslash}p{32mm}@{}}
    \textbf{#1} & \textcolor{Muted}{#3}\\[-0.1em]
    \textcolor{Teal}{#2} &
  \end{tabularx}
}
\newcommand{\pubtitle}[1]{\textbf{\color{Ink}#1}}
\newcommand{\meta}[1]{{\small\color{Muted}#1}}
\newcommand{\cvitem}[1]{\item {#1}}
\newenvironment{cvitems}{%
  \begin{itemize}[leftmargin=1.15em,itemsep=0.17em,topsep=0.25em,parsep=0pt,partopsep=0pt,label=\textcolor{Teal}{\raisebox{0.15ex}{\scriptsize\textbullet}}]
}{\end{itemize}}
\newcommand{\tag}[1]{\colorbox{Soft}{\strut\hspace{0.35em}\textcolor{Navy}{\small\bfseries #1}\hspace{0.35em}}}

\begin{document}

\noindent\color{Navy}\rule{\textwidth}{2.2pt}
\vspace{0.7em}

{\fontsize{27}{30}\selectfont\bfseries\color{Ink}Jia-Hao (Jay) Ji}\par
\vspace{0.18em}
{\large\color{Navy}\bfseries M.D. Candidate \textbar{} AI Researcher \textbar{} Physician-Scientist in Training}\par
\vspace{0.55em}
{\small\color{Muted}
Shanghai, China
\contactsep \href{mailto:jhji19@fudan.edu.cn}{jhji19@fudan.edu.cn}
\contactsep \href{https://jiahaoblog.com}{jiahaoblog.com}
\contactsep \href{https://github.com/Biogod2020}{github.com/Biogod2020}
\contactsep \href{https://orcid.org/0009-0007-4647-522X}{ORCID}
}\par

\vspace{0.85em}
\begin{minipage}{\textwidth}
\colorbox{Soft}{%
  \parbox{\dimexpr\textwidth-2\fboxsep\relax}{%
    \vspace{0.35em}
    \small\color{Ink}
    Eight-year M.D. candidate at Shanghai Medical College, Fudan University developing reliable agentic and multimodal AI systems for biomedical discovery. My research connects large-language-model agents, spatial omics, digital pathology, and neuroscience, with emphasis on measurable reliability and scalable scientific workflows.
    \vspace{0.35em}
  }%
}
\end{minipage}

\vspace{0.7em}
\tag{Agentic AI for Science}\hspace{0.45em}
\tag{Spatial Omics}\hspace{0.45em}
\tag{Multimodal Learning}\hspace{0.45em}
\tag{Multi-Agent Reasoning}\hspace{0.45em}
\tag{Neuroscience}

\section*{Selected Publications and Manuscripts}

\pubtitle{Candidate supply and answer selection shape the value of LLM judging in multi-agent systems}\par
\meta{\textbf{Jia-Hao Ji}, Sijie Li, Jiabei Cheng, Zixi She, Jin-Tai Yu, and Zhiyuan Yuan. arXiv:2608.25937 [cs.AI, cs.MA], 2026. \href{https://arxiv.org/abs/2608.25937}{Paper}}
\begin{cvitems}
  \cvitem{Separated candidate generation, answer recognition, and terminal selection in controlled multi-agent reasoning experiments across five benchmarks.}
  \cvitem{Analysed 15,336 questions and replayed 81,390 fixed candidate pools; changing only the final selection rule increased accuracy from 63.82\% to 70.82--70.95\%.}
\end{cvitems}

\vspace{0.28em}
\pubtitle{SpatialDataAgent: Autonomous Spatial Omics Data Curation at Decade Scale}\par
\meta{\textbf{Jia-Hao Ji}, Qi Zou, Jiabei Cheng, Zixi She, Yining Hao, Weishi Liu, Daoliang Zhang, Zhikang Wang, Jin-Tai Yu, and Zhiyuan Yuan. bioRxiv, 2026. \href{https://doi.org/10.64898/2026.05.27.727615}{DOI}}
\begin{cvitems}
  \cvitem{Conceived and led development of an evidence-grounded agentic workflow combining schema-constrained evaluation with self-refining data standardisation.}
  \cvitem{Recovered 769 paired H\&E--spatial transcriptomics datasets, expanded existing curated baselines by 6.4-fold, and assembled HESRT with 29.2 million spots or cells.}
\end{cvitems}

\vspace{0.28em}
\pubtitle{Spatiotemporally resolved transcriptome unveils reactive subpallial microglia underlying cholinergic vulnerability in Alzheimer's disease}\par
\meta{\textbf{Jia-Hao Ji} (co-first author) et al. Under review at \textit{Cell}, 2026.}
\begin{cvitems}
  \cvitem{Integrated spatial transcriptomics, single-nucleus sequencing, and in situ amyloid pathology across 12.36 million spatial profiles.}
  \cvitem{Characterised disease-associated microglial microenvironments and their relationship to basal forebrain cholinergic vulnerability.}
\end{cvitems}

\section*{Research Experience}

\roleline{Researcher and Project Lead}{Fudan Data-Driven Future Lab \& Institute of Medical Genetics, Fudan University}{Sep 2023 -- Present}
\meta{Advisors: Prof. Jin-Tai Yu and Prof. Zhiyuan Yuan \textbar{} Alzheimer's disease, spatial omics, and biomedical AI}
\begin{cvitems}
  \cvitem{Lead computational projects spanning large-scale spatial-data curation, multimodal tissue registration, and disease-focused analysis of histology and molecular profiles.}
  \cvitem{Developed transformer-based and image-registration pipelines to align whole-slide pathology with spatial transcriptomic coordinates and map beta-amyloid plaque microenvironments.}
  \cvitem{Designed reproducible agentic workflows with explicit schemas, evidence tracking, failure auditing, and multi-stage validation.}
\end{cvitems}

\vspace{0.3em}
\roleline{Undergraduate Researcher}{Prof. Bin Song's Lab, Fudan University}{Oct 2021 -- Jun 2023}
\meta{Parkinson's disease and single-cell transcriptomics}
\begin{cvitems}
  \cvitem{Built scRNA-seq analysis pipelines to identify markers relevant to the purity, safety, and therapeutic potential of autologous iPSC-derived dopaminergic progenitor transplantation.}
\end{cvitems}

\vspace{0.24em}
\roleline{Undergraduate Researcher}{Prof. Fei Lan's Lab, Institutes of Biomedical Sciences, Fudan University}{Sep 2019 -- Jul 2021}
\meta{DNA methylation, epigenetics, and molecular biology}
\begin{cvitems}
  \cvitem{Received wet-lab training in cell culture, molecular assays, and experimental analysis of DNA methylation and epigenetic regulation.}
\end{cvitems}

\newpage
\noindent\color{Navy}\rule{\textwidth}{2.2pt}
\vspace{0.55em}
{\small\bfseries\color{Navy}JIA-HAO JI}\hfill{\small\color{Muted}BIOMEDICAL AI \textbar{} SPATIAL OMICS \textbar{} MULTI-AGENT REASONING}\par

\section*{Education and Clinical Training}

\roleline{Doctor of Medicine, Eight-Year Clinical Medicine Program}{Shanghai Medical College, Fudan University}{Sep 2019 -- Jun 2027 (expected)}
\begin{cvitems}
  \cvitem{Clinical focus in neurology; research focus in biomedical AI, computational biology, spatial omics, and neurodegenerative disease.}
\end{cvitems}

\vspace{0.28em}
\roleline{Exchange Program}{University of California, Los Angeles}{Sep 2021 -- Feb 2022}
\meta{Interdisciplinary coursework and research-oriented academic training.}

\vspace{0.22em}
\roleline{Clinical Elective, Department of Medicine}{HKU Queen Mary Hospital}{Apr 2025 -- May 2025}
\meta{Ward rounds, case discussions, and multidisciplinary clinical teaching.}

\section*{Selected Systems and Engineering}

\roleline{ASSA -- Autonomous Self-Sovereign Agent}{Independent open-source project}{2026}
\begin{cvitems}
  \cvitem{Developed an agent framework with hierarchical memory, reflection, and skill-generation mechanisms for long-horizon adaptation. \href{https://github.com/Biogod2020/ASSA}{GitHub}}
\end{cvitems}

\vspace{0.28em}
\roleline{Spatial-CLIP}{Multimodal representation learning for spatial biology}{2025 -- Present}
\begin{cvitems}
  \cvitem{Fine-tuned vision-language encoders with contrastive objectives to align histopathology images and spatial transcriptomic states for cross-modal retrieval and analysis. \href{https://github.com/Biogod2020/Spatial-Clip}{GitHub}}
\end{cvitems}

\section*{Teaching and Scientific Communication}

\begin{cvitems}
  \cvitem{Delivered invited sessions for graduate and undergraduate students on practical AI workflows for literature synthesis, scientific writing, and reproducible research, Fudan University, December 2025.}
  \cvitem{Developed a course module on artificial intelligence in scientific discovery and biomedical research; maintain a technical research blog at \href{https://jiahaoblog.com}{jiahaoblog.com}.}
\end{cvitems}

\section*{Technical Expertise}

\begin{tabularx}{\textwidth}{@{}>{\bfseries\color{Navy}}p{35mm}X@{}}
Machine Learning & PyTorch, transformers, graph neural networks, contrastive learning, multimodal foundation models, LLM agents, judge-based evaluation, multi-agent systems \\
Biomedical Data & Spatial transcriptomics, single-cell RNA-seq, digital pathology, whole-slide imaging, tissue registration, statistical analysis \\
Engineering & Python, R, TypeScript, Linux, Git, multi-GPU training, high-performance computing, LangGraph, ReAct-style orchestration \\
\end{tabularx}

\vspace{0.55em}
{\large\bfseries\color{Navy}Honours and Languages}\par
\vspace{-0.35em}\color{Rule}\rule{\textwidth}{0.45pt}\par
\vspace{0.32em}
\color{Ink}
\begin{tabularx}{\textwidth}{@{}>{\bfseries\color{Navy}}p{35mm}X@{}}
Selected Honour & First Prize, National Biology Olympiad (China); ranked 9th in Guangdong Province \\
Languages & Mandarin, Cantonese, and Teochew (native); English (TOEFL 105); Spanish (basic) \\
\end{tabularx}

\end{document}
